\SourceLeaf{205}\begin{center}
\begin{tabular}{r|p{.82\linewidth}}
I & \\
 & Umgestellt\\
II & \\
 & An 7 geschraubt\\
III & \\
 & An 7 halb zurückgeschraubt\\
IV & \\
 & An 3 geschraubt, wenn zugleich die Libelle berichtigt werden soll.\\
V &
\end{tabular}
\end{center}

Die Schl\"usse r\"uckw\"arts sind also:

(1) $v^4=0$, weil (2) $v^3=-v^2$ und kleine gleiche Bewegungen von 7 das $v$ gleich viel ver\"andern.

(2) $v^3=-v^2$, weil $v^2=u^1$, $v^3=-u^3$ und (3) $u^1=u^3$.

(3) $u^1=u^3$, weil $S^1=S^3=u$ von $u$ abh\"angig, also $u^1=u^3$.

Da also $S^1$ hier in Consideration kommt, so ist klar, dass man vor I schon die Ebene $T$ in die Lage gebracht haben muss, dass die Enden der Blase sichtbar sind.

Soll nun noch die Berichtigung der Libelle bewirkt werden, so ist zu zeigen, dass $i^5=0$.

Diess erkennt man an der Normalstellung von $S$, aber offenbar h\"angt sie nicht unmittelbar mit $i=0$, sondern nur mit $g=0$ zusammen.

Die Synthesis geht nun am besten von den urspr\"unglichen Werthen aus:
\[
S = M+h+i+u,
\]
\[
\Sigma = M-h+i+u,
\]
\[
\begin{array}{rccc}
 & u & v & i \\
\cline{2-4}
{[\mathrm{I}]} & v^1 & v^1 & i^1 \\
{[\mathrm{II}]} & -v^1 & v^1 & i^1 \\
{[\mathrm{III}]} & v^1 & -v^1 & i^1 \\
{[\mathrm{IV}]} & 0 & 0 & i^1 \\
{[\mathrm{V}]} & 0 & 0 & 0 \\
\end{array}
\]

Da $v$ bei 5 verschiedenen Stellungen nur 1\,mal $=+u$ und 4\,mal $=-u$ ist, so kann man $v$ als Zeichen ganz ignoriren und [nur] von $u$ sprechen.

